\documentclass[11pt]{article}
\usepackage{amsmath,amssymb,hyperref}
\title{Recover a usable link while the formation is still braking}
\begin{document}\maketitle
\paragraph{Target and omissions.}
After a sudden stop, preserve clearance and every established
surviving old link throughout. Restore one cross-component link at
the activation range for the last \(h\) time units before a common
deadline \(T\); all surviving robots must be at rest at \(T\).
The acquisition interval is a geometric proxy for a transport
that can run in motion. No packet, throughput, control-message
delivery, or measured radio law is claimed. A protocol that
requires the sender and receiver to be stationary throughout its
transfer has a different task and cannot use the overlap. The
stopped failed body remains an obstacle. The fault can remove global
connectivity immediately, so only surviving old links are kept.

\paragraph{Exact scalar envelope.}
Normalize relative progress to \(x(0)=0\) and \(v(0)=v(T)=0\),
with \(\dot x=v\), \(|\dot v|\le A\), \(|v|\le V\) and
\(x(t)\le L\) as the clearance constraint, where \(A,V>0\).
Write \(C\ge0\) for the progress threshold activating the link.
Reverse motions are permitted in the necessity claim. With
\[
 F(t;\tau)=\int_0^{\min(t,\tau)}
  \min\{As,V,A(\tau-s)\}\,ds,\qquad I(\tau)=F(\tau;\tau),
\]
take \(\tau=\min(T,I^{-1}(L))\) for \(L\ge0\). Then for all
\(t\in[0,T]\), \(F(t;\tau)\) is the largest attainable position
with all these requirements. If \(I(T)\le L\), integrate the three
pointwise velocity limits from speed, initial rest and final rest.
Otherwise the profile saturates the clearance cap. Before its
braking switch, the corresponding speed envelope bounds every
trajectory. During braking its squared speed is exactly
\(2A(L-F)\). Any trajectory with positive current velocity needs at
least \(v^2/(2A)\) further progress before it can stop, so has
\(v^2\le 2A(L-x)\). Where a proposed trajectory were ahead of
the profile during braking, it would thus have strictly lower
speed (negative speeds cannot overtake either). The positive
part of the gap to the profile is absolutely continuous and
cannot increase there; with no initial positive gap it stays
zero. After the profile stops, the clearance cap itself is tight.
The profile is realizable by accelerating, cruising if necessary,
braking and then holding. Therefore the joint demand is exactly
\[
 T\ge h,\qquad C\le L,\qquad F(T-h;\tau)\ge C .
\]
The monotone realizing profile remains linked after the left end
of the dwell. A two-component symmetric axial translation
realizes the relative bound with twice the individual acceleration
and speed. That lift assumes common timing and no tighter
intermediate constraint than the stated axial clearance.

\paragraph{A joint incompatibility despite separate feasibility.}
Take \(A=1,V=4,C=1/4,L=3/5,h=2/5\). Rest-only needs time 1,
while acquiring the link for the full dwell without stopping at
\(T=28/25\) is possible: accelerate until \(71/100\), then brake,
giving sufficient progress by \(18/25\) and positive velocity at
the deadline. For joint rest and dwell the largest possible
position at \(T-h\) is \(T^2/4-h^2/2<1/4\). The exact joint
deadline is \(\sqrt{33/25}\), instead of 1.4 if no progress
toward the dwell is admitted before rest. With tight clearance
\(L=C\), first attaining the link necessarily reaches the cap
at rest, and the 1.4 deadline is exact. The replacement of
rest-first evidence by moving evidence works only where an
extra safe closing distance is actually available.
Lean checks the braking-curve comparison and rational
incompatibility witness; it does not formalize the full ODE.

\paragraph{Native warm states and universal lower bound.}
Re-run to the four request cuts of the two published
\href{https://arxiv.org/html/2609.29939v1}{KTH managed-removal point-mass examples}
from the pinned source. At each departure request, counterfactually
park the outgoing body and drop its edges while preserving both
positions and velocities of surviving bodies. The source radii
and old graph are reused. Speed cap 1, acceleration cap \(A\)
and dwell duration are imposed in the recovery test, not measured
vehicle capabilities. We impose no disturbance or state error.
For any potential bridge at time \(T-h\), project both velocities
towards the other along the initially separating axis.
A robot with initial signed projection \(u\) can close by
at most \(\int_0^{T-h}\min\{V,u+At,A(T-t)\}\,dt\).
The sum must cover at least the initial separation minus range,
for at least one pair. This is necessary even if pair identity
could change after the start of the dwell, and even if all other
constraints were removed. Rest at \(T\) also needs at least the
longest individual braking time.

\paragraph{Constructive upper bound and checks.}
Brake each survivor from its actual velocity, then move both
components symmetrically along the closest braked cross pair.
Choose a closer rest endpoint than the old activation radius
so that the crossing occurs before rest. The same identified
bridge distance decreases linearly in the axial displacement
during the entire dwell. Piecewise constant acceleration and
a chord-deviation bound check every old edge, body and obstacle
continuously, not just at sample endpoints. Number arithmetic on
the native inputs is floating point with strict witness and
rejection margins, not exact real root certification for arbitrary
data. All 36 pose/cap/dwell combinations give witnesses within
0.00101 model units of the relaxed motion lower bound and save at
least 0.1054 relative to the lower bound under the rest-first
law. For the planar request at acceleration 1 and dwell 0.25,
the bound is about 1.02193 and the witness 1.02244, versus at
least 1.20879 under rest-first. The ending pair gap is about
1.47749, below the service radius 1.54; old links and clearance
have their own measured certificates. The replay also cross-checks
3,348 scalar-profile evaluations.

\paragraph{Baselines, interface and residual.}
Managed removal establishes a bridge before permitted departure;
that is a way to avoid the abrupt starting condition, not a
matched controller after an unannounced stop. The
\href{https://en.endo-lab.jp/research/detail/499}{Maeda--Endo--Matsuno
work and later distributed MPC} are substantive local
crash-recovery comparisons. Their detailed controllers and
hardware have not been reproduced. Breeden and Panagou's
\href{https://arxiv.org/abs/2210.01354}{joint input viability}
protects invariant domains, whereas here time to acquisition
and terminal rest also matter. A same-plant direct solver with
the more permissive mobile seal can exploit the same trajectories;
the reported advance is an exact geometric boundary and a
measured effect of changing the service interface, not an
advantage over that solver. Warm replay uses globally known
state at the cut, not a local-information implementation.
Whether a real acquisition works in motion, is long enough,
and can confirm the correct application or authority receipt is
unresolved. No independent mission or full R6 claim is made.
\end{document}
